\section{Discussion}\label{sec:discussion}

\subsection{What the structure does and does not do}

The results of Sections~\ref{sec:kernel}--\ref{sec:cm} are kinematics on one point.  They show that the hypothesis of one object in three
frames has a precise and rather rigid form on the kernel: the Lorentz invariant \(s\) separates the sectors, duality exchanges the massive
ones, light is the duality-invariant fix-sector, and the symmetric triple is the centred null configuration.  None of this assigns a field
equation to \(z\).  The compensator is a projection, not a law of motion, and on the doublet the operators compatible with it only weight its
two ports.

Section~\ref{sec:generation} reaches further, but under a stated setting: the weak \(SU(2)\) acting from the right of \(\HH\), the form of the
electric charge, and a chirality projector acting on \(\mathbb C\otimes\OO\).  Within that setting nothing is left to choose.  Why the setting
holds is not explained, nor are symmetry breaking, the values of the masses, or the mixing angles.

The numbers of Section~\ref{sec:pheno} inherit the status of \(\beta_{\rm sym}\), which is put in.  The counting reading makes it a ratio of two
countings rather than a fitted parameter, and it is consistent with an independent window from lepton mixing, but it selects the U(1)
treatment of the horizon over the SU(2) treatment that the literature prefers.  An origin of \(\beta_{\rm sym}\) that does not depend on this
choice would change the status of every readout in that section.

\subsection{What a dynamics would have to do}

The common gap is the absence of an action varied in all fields.  The results above constrain such an action fairly precisely.
\begin{enumerate}
\item Coframe, connection and matter must be independent fields that act back; a geometry assembled from probe fields alone becomes flat.
\item Matter and tachyon need the same kinetic term, so that the quarter turn \eqref{eq:quarterturn} is a symmetry of the action.
\item The compensator must enter as a coupling, not as a projector applied afterwards.
\item \(\Lambda\) should appear as an integration constant, as in unimodular gravity~\cite{henneaux} or vacuum-energy sequestering~\cite{kaloper}:
state, frame and normalization never fix the absolute vacuum term.
\item A boundary term that counts one \(U(1)\) --- the inner unit \(L_{e_7}\) selects one \(U(1)\) in \(SU(2)\) --- might ground the counting rule
behind \(\beta_{\rm sym}\).
\item Triality should be a symmetry of the kinetic term, broken by chirality.
\item A clock should come from the phase difference between matter and tachyon, since a common energy offset is not observable.
\end{enumerate}
A first test shows a limit of this programme: in the simplest reduction a frozen dimensionful value of \(\beta_{\rm sym}\) always carries the
expansion history.  A dimensionless structural value has to come from a dimensionless ratio, such as the two countings above.

\subsection{Field equations}

A comparison model of the probe field equations, in which the coframe is reconstructed from the spinor field itself, has been studied with
the formal theory of partial differential equations~\cite{bryant,seiler}.  It has an exact gauge symmetry; after gauge fixing its symbol
becomes involutive at order six, with solutions depending formally on \(14\) functions of three variables; and four hidden conditions of
sixth order appear, which act only on light-like plane waves.  Whether the system extended by them is formally integrable is open.  These
results concern a declared model, not a derivation of the field equations from the kernel, and they will be reported separately.

\subsection{Open questions}

Several questions can be studied without accepting any physical reading of the framework.  Is the anomaly selection of
Proposition~\ref{prop:anomaly} known, and does it survive other placements of the weak \(SU(2)\)?  Is there a realization of the
compensator-matrix conditions on a 16-dimensional carrier, for instance one generation, that satisfies the hypotheses of
Proposition~\ref{prop:spectral}?  Is there a horizon counting in the SU(2) treatment whose minimal and full versions differ by the factor
\(1.8647\)?  And is the extended comparison system formally integrable?
